% Глава 7: Заключение / Chapter 7: Conclusion

\chapter{Заключение}
\label{ch:conclusion}

\section{Основные результаты}
\label{sec:concl:results}

В диссертации решена актуальная научная задача автоматизации ввода и верификации больших данных в интерактивных наукометрических системах. Получены следующие основные результаты:

\begin{enumerate}
    \item \textbf{Разработан мультимодульный алгоритм} автоматической верификации порядка имени и фамилии в китайских именах авторов, интегрирующий лингвистические, статистические и контекстные признаки.

    \item \textbf{Создана архитектура с цепочками доказательств}, обеспечивающая полную интерпретируемость каждого решения системы и возможность аудита результатов.

    \item \textbf{Реализован механизм автоматической редакции персональных данных} на основе хеширования SHA-256 с солью, соответствующий принципам GDPR и ФЗ-152.

    \item \textbf{Проведена экспериментальная оценка} на крупномасштабных реальных данных Crossref (N=301~446), показавшая точность [TODO: insert value]~\% и пропускную способность [TODO: insert value]~имен в секунду.

    \item \textbf{Выполнено внедрение} разработанного решения в промышленную эксплуатацию системы ИСТИНА МГУ, обрабатывающей более 10~000 китайских имен авторов в рамках регулярного импорта данных.
\end{enumerate}

\section{Вклад в науку}
\label{sec:concl:contribution}

Научная новизна работы заключается в следующем:

\begin{itemize}
    \item Предложен новый подход к автоматической верификации данных, основанный на мультимодульной архитектуре с цепочками доказательств, что обеспечивает превосходную интерпретируемость по сравнению с существующими методами.

    \item Впервые создан крупномасштабный размеченный датасет для задачи верификации порядка имени и фамилии в китайских именах на основе данных Crossref и ORCID (N=301~446).

    \item Экспериментально подтверждена эффективность комбинирования лингвистических, статистических и контекстных признаков для решения задачи верификации многоязычных имен.
\end{itemize}

\section{Практическая значимость}
\label{sec:concl:practical}

Практическая ценность работы подтверждается:

\begin{itemize}
    \item Внедрением разработанного решения в промышленную эксплуатацию системы ИСТИНА МГУ, используемой более чем [TODO: insert number] исследователями и администраторами.

    \item Повышением качества и скорости обработки библиографических данных в системе ИСТИНА, что способствует более точной оценке научной деятельности.

    \item Возможностью адаптации предложенной архитектуры для решения других задач верификации данных в интерактивных системах (дедупликация авторов, нормализация аффилиаций и т.д.).

    \item Открытым доступом к исходному коду и воспроизводимым результатам, что способствует развитию исследований в области автоматизации обработки научных данных.
\end{itemize}

\section{Направления дальнейших исследований}
\label{sec:concl:future}

Перспективные направления развития работы:

\begin{enumerate}
    \item \textbf{Расширение на другие языки}: Адаптация предложенного подхода для обработки корейских, вьетнамских, арабских и других многоязычных имен.

    \item \textbf{Интеграция машинного обучения}: Использование методов обучения с подкреплением (reinforcement learning) для автоматической оптимизации весов модулей на основе обратной связи от экспертов.

    \item \textbf{Применение в смежных задачах}: Адаптация архитектуры с цепочками доказательств для задач дедупликации авторов, нормализации аффилиаций и извлечения знаний из научных текстов.

    \item \textbf{Федеративное обучение}: Разработка методов распределенного обучения и верификации данных в интерактивных системах с сохранением конфиденциальности.
\end{enumerate}

\section{Заключительные замечания}
\label{sec:concl:final}

Данная работа вносит вклад в решение актуальной проблемы автоматизации ввода и верификации больших данных в интерактивных наукометрических системах. Предложенный мультимодульный подход с цепочками доказательств обеспечивает баланс между точностью, производительностью и интерпретируемостью, что критично для промышленных применений.

Внедрение результатов работы в систему ИСТИНА МГУ подтверждает их практическую значимость и готовность к реальной эксплуатации. Открытый доступ к коду и данным способствует воспроизводимости исследований и дальнейшему развитию методов автоматизации обработки научных данных.
